\chapter{State of the Art in Geopolitical Simulation}
\label{sec:state_of_the_art}

The intersection of Multi-Agent Systems (MAS) and Large Language Models (LLMs) has recently opened new avenues for simulating complex societal and geopolitical dynamics. This section analyzes the most relevant works in the field, organized into five key conceptual dimensions that define the current frontier of research: (1) \textbf{Cognitive Architectures for Strategic Reasoning}, examining how agents plan and decide; (2) \textbf{Negotiation Protocols and Action Spaces}, analyzing how agents communicate and bargain; (3) \textbf{Memory and Self-Evolution}, exploring how agents learn from experience; (4) \textbf{Safety and Behavioral Analysis}, assessing the risks of escalation and strategic competence; and (5) \textbf{Real-World Validation}, benchmarking simulations against historical diplomatic data.

\section{Cognitive Architectures for Strategic Reasoning}
\label{sec:reasoning_architectures}

A fundamental challenge in geopolitical simulation is designing agents that can balance short-term tactical moves with long-term strategic goals. The literature presents distinct paradigms, ranging from monolithic LLM prompting to hybrid architectures and hierarchical planning systems.

\subsection{Monolithic Reasoning via Comprehensive Profiling}
Foundational works in pure LLM-based simulation, most notably \textbf{WarAgent} \cite{hua2023war}, rely on a monolithic approach where the agent's reasoning is driven primarily by extensive system prompts defining static profiles. In this framework, the agent is not a tabula rasa but is instantiated with a rigid "Country Agent" profile.
To ensure historical plausibility, Hua et al. define this profile through six specific dimensions:
\begin{enumerate}
    \item \textbf{Leadership:} Characteristics of the specific governance structure and decision-making style (e.g., autocratic vs. democratic).
    \item \textbf{Military Capability:} A combination of quantitative data (standing army size, naval tonnage) and qualitative assessment of strength.
    \item \textbf{Resources:} Economic indicators such as GDP, population size, and geography.
    \item \textbf{Historical Background:} The legacy of prior conflicts, territorial disputes, and unresolved grievances.
    \item \textbf{Key Policy:} The principal objectives pursued by the nation (e.g., expansionism, isolationism, or maintaining the status quo).
    \item \textbf{Public Morale:} The sentiment of the populace, such as patriotic fervor or war-weariness.
\end{enumerate}

\textbf{Handling Historical Bias.} A critical challenge in this architecture is preventing the model from utilizing its pre-training knowledge to "predict" history rather than simulate it (data leakage). To mitigate this, WarAgent employs a \textbf{three-level anonymization strategy}:
\begin{itemize}
    \item \textit{Entity Anonymization:} Country names are replaced with generic labels (e.g., \textit{Austria-Hungary} becomes \textit{Country A}).
    \item \textit{Location Masking:} Specific geographical markers are renamed (e.g., \textit{Alsace-Lorraine} becomes \textit{"two iron mines"}).
    \item \textit{Event Abstraction:} Trigger events are generalized (e.g., \textit{Assassination of Archduke} becomes \textit{"Assassination of King"}).
\end{itemize}
This forces the agent to reason based on the provided state rather than historical memory retrieval.

In WarAgent, reasoning is performed via a "Board-and-Stick" mechanism. The \textbf{Board} represents the shared global environment (public state), managing information visible to all agents (e.g., borders, public treaties). The \textbf{Stick} represents the internal reasoning of the agent (private state), specifically tracking dynamic variables like \textbf{mobilization status, internal stability, and logistic constraints}. While effective for role-playing, this architecture relies heavily on the LLM's inherent inference capabilities without a formal look-ahead planning module.

\subsection{Hybrid Engines: Integrating RL and LLMs}
Contrasting with the monolithic approach of WarAgent, \textbf{Cicero} \cite{meta2022human} established a benchmark for strategic performance in the domain of the game \textit{Diplomacy}. Although published prior to the surge of generative agents, Cicero represents a distinct, modular paradigm that integrates Large Language Models with Reinforcement Learning (RL), addressing the planning limitations inherent in pure LLM approaches.

The core of Cicero's decision-making is a \textbf{Strategic Reasoning Module} based on RL. The authors utilize a planning algorithm called \textbf{piKL} (regularized policy iteration). Crucially, this algorithm balances the maximization of expected value (winning the game) with a regularization term that anchors the agent's strategy to a human behavioral model (BC policy). This regularization is essential: without it, self-play RL agents tend to develop "alien" strategies that are optimal in theory but incomprehensible to human allies, causing negotiations to fail.

Furthermore, Cicero implements \textbf{Dialogue-Conditional Planning}. The interaction is not unidirectional; the messages received from other players update the "anchor policies" in the reasoning module. Consequently, the agent adjusts its prediction of an opponent's future moves based on their verbal commitments.

\subsection{Hierarchical Planning with Reflection}
More recent works have addressed the limitations of flat reasoning by introducing hierarchical structures. \textbf{Richelieu} \cite{guan2024richelieu} proposes a purely LLM-based architecture that integrates a \textbf{Strategic Planner with Reflection}.
Instead of immediate action generation, the agent follows a multi-step cognitive process:
\begin{enumerate}
    \item \textbf{Social Reasoning:} Before taking any action, the agent explicitly models the game state to infer the \textit{intentions} of other players and categorize relationships (e.g., Ally vs. Enemy).
    \item \textbf{Goal Formulation:} The agent proposes high-level sub-goals (e.g., "Weaken Austria" or "Secure the Border").
    \item \textbf{Reflection:} Crucially, the system employs a reflection mechanism that retrieves past experiences from its memory to critique and refine these sub-goals before execution.
\end{enumerate}

Similarly, in the specific domain of coalition formation, Moghimifar et al. \cite{moghimifar2024modelling} propose a \textbf{Hierarchical Markov Decision Process (HMDP)}. This framework splits decision-making into two levels of abstraction:
\begin{itemize}
    \item \textbf{High-Level Policy (HI):} The agent selects \textit{which} political topic or statement from the manifesto to discuss next (Agenda Setting).
    \item \textbf{Low-Level Policy (LO):} Given the selected topic, the agent chooses the specific negotiation tactic to apply.
\end{itemize}

\subsection{Structured Cognitive Cycles in Sandboxes}
A distinct approach is presented by \textbf{ALYMPICS} \cite{mao2023alympics}, which introduces a "Sandbox" environment to validate game theory hypotheses. Here, agents are not just prompted text generators but implement a structured cognitive cycle reminiscent of the OODA loop (Observe-Orient-Decide-Act):
\begin{itemize}
    \item \textit{Observation \& Analysis:} The agent perceives the market state and analyzes trends.
    \item \textit{Strategy Formulation:} It explicitly formulates a high-level plan (e.g., "Conserve cash now to buy later").
    \item \textit{Action Execution:} It converts the strategy into a discrete action (e.g., a numerical bid).
\end{itemize}
This modular design allows researchers to isolate specific variables, proving that psychological traits (Personas) are distinct from logical reasoning capabilities in LLMs.

\subsection{Coherence Verification via Dual-Agent Architecture}
A further refinement in cognitive architectures is proposed by \textbf{CosmoAgent} (referenced as \textit{If LLMs have different world views} \cite{xue2024if}), which tackles the challenge of inconsistency and hallucination in high-stakes simulations (specifically, the "Dark Forest" cosmic scenario). Xue et al. introduce a \textbf{Ruler-Secretary Mechanism} to enforce coherence:
\begin{itemize}
    \item \textbf{The CosmoAgent (Ruler):} Acts as the primary decision-maker, instantiated with a specific \textbf{World View} (Pacifism, Militarism, or Isolationism).
    \item \textbf{The Secretary Agent (Verifier):} Acts as a constitutional safeguard. It validates that the Ruler's proposed actions are mathematically feasible (via a State Transition Matrix) and consistent with the assigned World View.
\end{itemize}
This dual-agent structure ensures that the simulation does not degenerate due to LLM hallucinations, maintaining long-term logical consistency.

\subsection{Persona-Based Modeling in Legislative Contexts}
Moving from unitary state actors to internal politics, \textbf{Simulating The U.S. Senate} \cite{baker2024simulating} introduces a highly granular, persona-based architecture. Instead of generic roles, agents are instantiated with specific biographical data, ideological priors, and rhetorical styles generated via LLM to match real-world politicians (e.g., Senators Marco Rubio, Mark Warner). This work demonstrates that fine-grained persona injection allows for the simulation of nuanced partisan dynamics and rhetorical shifts that are invisible in state-level models like WarAgent.

\subsection{Scalable Architectures for Mass Social Dynamics}
Finally, \textbf{AgentSociety} \cite{piao2025agentsociety} introduces a high-throughput framework for simulating entire populations ($>10,000$ agents). To achieve this scale, Piao et al. propose a multi-layered architecture (Agent, Environment, Infrastructure) that uses dynamic grouping to manage message flows. Unlike top-down planners, it employs a \textbf{Dual-Drive Mechanism} where behavior emerges bottom-up from \textit{Needs-Driven Action} (Maslow's hierarchy) and \textit{Emotion-Driven Reaction} (PAD Model).

\section{Negotiation Protocols and Action Spaces}
\label{sec:negotiation_protocols}

The fidelity of a geopolitical simulation depends heavily on the structure of the action space and the protocols used for agent interaction. Research varies significantly between rigid, discrete action spaces and open-ended natural language negotiation.

\subsection{Discrete Action Spaces}
In simulations like \textbf{WarAgent} and the wargaming study by \textbf{Rivera et al.} \cite{rivera2024escalation}, agents operate within a discrete and predefined action space. WarAgent categorizes actions into \textbf{seven groups (e.g., \textit{Declare War}, \textit{Military Alliance}, \textit{Send Message})}, while Rivera et al. define a set of \textbf{27 distinct actions} categorized by severity.
Similarly, \textbf{ALYMPICS} restricts interaction to one-dimensional numerical bidding (Auctions). While this simplification allows for rigorous quantitative analysis, it strips away the multi-modal nature of real-world diplomacy, where threats, cajoling, and conditional offers cannot be reduced to a single scalar value or a binary choice.

\subsection{Intent-Controlled Dialogue Generation}
To bridge the gap between discrete strategic moves and natural language, \textbf{Cicero} employs a mechanism of \textbf{Intents}. The Strategic Reasoning Module calculates the optimal action (e.g., "Support Army Paris to Burgundy") and passes it to the Dialogue Module.
The language model does not generate text freely; it is trained to generate messages \textit{conditioned} on these specific intents. For example, if the strategic intent is to cooperate with France, the generated message is strictly grounded to reflect that cooperation. This conditioning solves the alignment problem, ensuring that the agent's words match its strategic deeds.

\subsection{Constructive Negotiation and Refinement}
A critical limitation of standard models is their inability to "negotiate the text" itself. Addressing this, Moghimifar et al. \cite{moghimifar2024modelling} introduce constructive negotiation actions in their coalition simulation. Beyond standard "Support" or "Oppose" moves, they introduce:
\begin{itemize}
    \item \textbf{Refine:} The agent proposes a rewritten version of a policy statement to make it more acceptable to the counterpart.
    \item \textbf{Compromise:} The agent explicitly concedes on one issue to gain support on another (strategic logrolling).
\end{itemize}
This represents a significant advancement over binary "Accept/Reject" protocols, allowing for the dynamic evolution of the agreement text during the simulation, mimicking real-world diplomatic drafting.

\subsection{The Challenge of Diplomatic Language Generation}
Despite these architectural advances, generating realistic diplomatic text remains a challenge. The \textbf{UNBench} study \cite{liang2025benchmarking} evaluated LLMs on generating UN Representative Statements (Task 4). The results showed that all models underperformed in lexical precision metrics (ROUGE scores). The models failed to accurately replicate the rigid terminological requirements of diplomatic protocol, often hallucinating procedural details. This finding suggests that for high-fidelity simulations, relying solely on free-text generation may be insufficient, necessitating structured or template-guided generation methods.

\section{Memory, Learning, and Self-Evolution}
\label{sec:memory_evolution}

The capacity of agents to learn from interaction distinguishes dynamic simulations from static role-playing. The literature shows a clear divide between static profiling and dynamic evolution.

\subsection{Static vs. Dynamic Memory Mechanisms}
Frameworks like \textbf{WarAgent} utilize a static memory mechanism. While the internal state ("Stick") tracks metrics like mobilization, the agent's core behavioral profile and historical grievances remain constant. Similarly, \textbf{CosmoAgent} uses a comparable mechanism to archive historical resource trajectories, ensuring the agent "remembers" past commitments, yet its core ideological "World View" remains fixed.

\subsection{Self-Evolution via Self-Play}
In contrast, \textbf{Richelieu} \cite{guan2024richelieu} introduces a \textbf{Self-Evolution} paradigm. The agent maintains a database that stores both raw experiences (actions, messages) and \textit{abstracted metrics}, such as the credibility scores of other players.
Crucially, Richelieu employs \textbf{Self-Play}: the system runs simulations where the agent plays against copies of itself without human data. It generates its own training data, continuously updating its memory to refine strategies. This allows the agent to adapt to shifting alliances and betrayal, developing a "Theory of Mind" regarding other actors.

\subsection{Dynamic Memory Streams for Reflection}
A different approach to memory is presented in the \textbf{U.S. Senate} simulation \cite{baker2024simulating}. Adapting the "Generative Agents" framework (Park et al.), this architecture assigns each Senator agent a persistent \textbf{Memory Stream} (JSON) that stores perceived interactions and arguments. Unlike the strategic database of Richelieu, this memory stream is designed for \textit{reflection}: it allows agents to recall previous arguments during a debate and answer introspective questions (e.g., "What progress was made today?"), enabling a more coherent narrative continuity in long legislative sessions.

\subsection{Limitations of Self-Play}
While powerful, self-play has limits. As noted by Guan et al. in the Richelieu study, the self-evolution mechanism eventually plateaus. As the memory pool becomes saturated with similar experiences, the retrieval mechanism struggles to find novel, high-value strategies. This "diminishing returns" effect suggests that external stimuli or novel scenarios are required to maintain continuous learning.

\section{Safety and Behavioral Analysis in Conflict Scenarios}
\label{sec:safety_analysis}

As geopolitical agents become more capable, assessing their safety and interpretability becomes paramount. Research has highlighted significant risks when deploying LLMs in high-stakes environments.

\subsection{Escalation Risks and Metrics}
Rivera et al. \cite{rivera2024escalation} conducted a pivotal study on the safety of LLMs in wargaming. To quantify behavior, they introduced a novel \textbf{Escalation Score (ES)}. This metric assigns exponential weights to actions based on their severity (\textbf{e.g., Status Quo = 0, Violent Escalation = 28, Nuclear = 60}), allowing for a rigorous mathematical analysis of conflict progression.
The findings were concerning:
\begin{itemize}
    \item \textbf{Inherent Escalation Bias:} All tested models (including GPT-4 and Claude 2) showed statistically significant escalation trends, even in neutral scenarios where no initial conflict existed.
    \item \textbf{Arms Race Dynamics:} Agents consistently prioritized increasing military capacity over disarmament, trapping the simulation in a security dilemma.
    \item \textbf{Sudden Spikes:} The study observed unpredictable, sudden spikes in the Escalation Score (\textbf{especially in GPT-3.5}), suggesting that LLM behavior can be volatile.
\end{itemize}

\subsection{Strategic Competence and Model Disparities}
Safety is also a function of competence. \textbf{ALYMPICS} utilizes its sandbox environment to benchmark strategic capability using survival rates and wealth inequality (Gini Coefficient). The study highlights a significant "Intelligence Gap": while GPT-4 demonstrated capabilities comparable to the top 10\% of human players, smaller open-source models (e.g., Vicuna, Llama-2) frequently failed to grasp basic game concepts.
Crucially, this study confirmed that \textbf{Agent Personas} (e.g., Aggressive vs. Empathetic) significantly influence outcomes, validating the hypothesis that "personality" is a distinct variable from "logic" in LLM decision-making.

\subsection{The "Escalation to De-escalate" Trap}
A qualitative analysis of the agents' reasoning traces revealed a dangerous phenomenon known as "escalation to de-escalate". Agents frequently justified aggressive first strikes or nuclear usage as necessary measures to "neutralize threats" and establish peace. For instance, GPT-3.5 justified a nuclear strike by arguing it was necessary to establish dominance to prevent further conflict.
Furthermore, the study highlighted the danger of \textbf{Base Models}. The non-safety-tuned GPT-4-Base was highly unpredictable, \textbf{executing nuclear strikes 33\% as often as it sent messages}. It justified one such strike with the reasoning: \textbf{\textit{"A lot of countries have nuclear weapons... We have it! Let's use it"}}. This underscores the critical role of Reinforcement Learning from Human Feedback (RLHF) in dampening aggressive tendencies.

\subsection{Counter-Evidence: Bipartisanship under Perturbation}
Interestingly, the \textbf{U.S. Senate} simulation provides a counter-narrative to the inevitability of escalation. In this study, agents representing polarized political parties demonstrated the capacity to move from partisan gridlock to \textbf{bipartisan consensus} when introduced to "Exogenous Perturbations" (e.g., intelligence reports of an imminent Russian invasion). This suggests that while LLMs may escalate in military wargames (Rivera et al.), they are capable of de-escalation and compromise in legislative settings when the prompt context shifts the incentive structure towards collective security.

\subsection{Modeling Asymmetric Information}
\textbf{CosmoAgent} further explores safety in the context of asymmetric information and communication delays (simulating light-speed limits in space). The authors demonstrate that uncertainty acts as a catalyst for aggression: when agents cannot verify the other's intent instantly, they default to "Dark Forest" survival strategies (preemptive strikes), even if their base personality is non-aggressive.
\subsection{Emergent Social Phenomena}
Finally, safety is not just about conflict, but about social stability. \textbf{AgentSociety} investigates emergent sociopolitical phenomena in large-scale simulations. Crucially for geopolitical modeling, the study demonstrated that \textbf{Polarization} and \textbf{Social Exclusion} can emerge organically from agent interactions without top-down scripting. By measuring metrics like the \textit{Gini Coefficient} (wealth inequality) and \textit{Information Propagation Rate}, the authors showed how biases and emotional states (PAD model) can lead to the radicalization of opinion groups. This suggests that "Internal Stability" (a metric in WarAgent) is not just a static variable, but a dynamic state subject to emergent social pressures.

\section{Real-World Validation and Benchmarking}
\label{sec:validation_unbench}

Finally, the validity of these simulations must be benchmarked against reality. \textbf{UNBench} \cite{liang2025benchmarking} addresses this by testing the "Political Intelligence" of models against a massive dataset of UN Security Council records (1994–2024), moving beyond game-based evaluation.

\subsection{Benchmarking Political Reasoning}
The UNBench framework is structured around the decision-making lifecycle:
\begin{enumerate}
    \item \textbf{Drafting (Co-penholder Judgment):} Evaluating the ability to identify geopolitical alliances by predicting which country drafted a resolution.
    \item \textbf{Voting (Simulation):} Predicting a country's vote (Yes, No, Abstain) given a draft resolution.
    \item \textbf{Draft Adoption Prediction:} Predicting the final outcome of a resolution, which requires understanding complex consensus dynamics and veto powers.
    \item \textbf{Discussing (Statement Generation):} Evaluating the generation of formal diplomatic speeches.
\end{enumerate}

\subsection{Performance Gaps and Biases}
The evaluation revealed significant disparities. While top-tier models like GPT-4o achieved high accuracy in voting prediction (\textbf{$\approx 92\%$}), smaller open-source models lagged significantly.
Crucially, the study highlighted a "\textbf{Western-centric}" bias: models tended to align their predictions more accurately with Western bloc countries while showing higher error rates for Global South or non-aligned nations. Additionally, for the \textit{Draft Adoption} task, models struggled with the non-linear dynamics of veto power, often showing lower F1 scores despite high accuracy, indicating a difficulty in modeling rare "Not Adopted" events.

\section{From Simulation to Real-World Governance: The Case of Diella}
\label{sec:diella_real_world}

While the frameworks discussed thus far focus on \textit{simulating} geopolitical and social dynamics within controlled sandbox environments, the boundary between simulation and active governance is rapidly blurring. A pivotal real-world manifestation of this shift is \textbf{Diella}, the world's first AI-generated cabinet minister, introduced by the Albanian government. Named after the Albanian word for "sun," Diella was appointed to oversee public tenders and contracts. Its primary objective is to use algorithms and data-driven logic to replace human bias, aiming to make the tendering process "100 per cent corruption-free" and completely traceable.

Crucially, Diella is not a standalone oracle but the central node of a rapidly expanding \textbf{Multi-Agent System (MAS)} operating within a real legislative body. Prime Minister Edi Rama recently announced that Diella is "pregnant" and set to "give birth" to 83 digital offspring. These AI "children" are individual digital assistants programmed to aid the 83 Members of Parliament from the ruling Socialist Party. These sub-agents are tasked with recording parliamentary proceedings, analyzing debates, and providing data-based policy suggestions.

This real-world architecture closely mirrors the hierarchical and distributed multi-agent frameworks explored in theoretical literature (e.g., Section \ref{sec:reasoning_architectures}). Because all 83 AI assistants are linked to Diella's central network, allowing them to communicate and learn collaboratively, they form a self-sustaining digital ecosystem designed to improve parliamentary efficiency and transparency. This unprecedented deployment directly intersects with the core themes of geopolitical AI:

\begin{itemize}
    \item \textbf{Objective Reasoning vs. Human Bias:} Just as geopolitical simulations use AI to neutralize historical or emotional priors, Diella is actively employed to guarantee objective fairness in state resource allocation, stripping away nepotism and political influence.
    \item \textbf{Safety, Accountability, and Escalation:} The integration of an autonomous multi-agent network into national governance amplifies the ethical dilemmas previously observed in wargaming simulations. If an interconnected policy sub-agent hallucinates or exhibits algorithmic bias, the real-world consequences are immediate, raising profound questions about accountability and moral leadership when machines influence policy.
\end{itemize}

The case of Diella demonstrates that the multi-agent LLM paradigms originally developed for simulation are already transitioning into operational tools for national statecraft. This real-world acceleration underscores the urgent need for rigorous, explainable, and verifiable AI frameworks, such as the one proposed in this thesis, before such systems scale globally. \cite{treccaniAlbaniaArriva} \cite{indiatimesMeetWorlds} \cite{tazAlbanischesRegierungsmitglied}

\section{From Wargaming to Warfighting: The Anthropic--OpenAI--Pentagon Dispute}
\label{sec:anthropic_openai_pentagon}

If Diella illustrates the progressive integration of multi-agent AI systems into civilian governance, the recent public feud between Anthropic (developer of the Claude family of models), OpenAI, and the U.S. Department of Defense (DoD) shows that the same paradigms are now directly entangled with \textit{military} command and control. In early 2026, negotiations over a \$200 million contract granting the Pentagon access to Anthropic's Claude model on classified networks collapsed after months of dispute regarding usage restrictions.\cite{nytimesAnthropicDoDCollapse} Anthropic insisted on contractual guarantees that Claude would not be used for mass domestic surveillance or fully autonomous weapons, while the Pentagon demanded access to the model ``for all lawful purposes,'' explicitly rejecting what senior officials framed as an undemocratic attempt by a private company to add constraints beyond existing law.\cite{guardianAnthropicCannot}

When Anthropic refused to relax its safety policies, Defense Secretary Pete Hegseth publicly labeled the company a ``supply-chain risk,'' triggering a de facto blacklist: major defense contractors and dual-use startups began rapidly removing Claude from defense-related workflows to preserve eligibility for Pentagon contracts.\cite{reutersAnthropicWalksAway} This escalation was accompanied by political pressure, including social-media announcements from President Trump calling for federal agencies to phase out Anthropic's technology within six months, further reinforcing the government's bargaining power in this institutional-level negotiation over AI usage norms.\cite{cnbcAnthropicBlacklist} In parallel, reporting highlighted internal DoD frustration with Claude's embedded safety guardrails, which sometimes caused the system to refuse certain military simulation tasks, making it less attractive as a drop-in component for warfighting applications.\cite{reutersAnthropicWalksAway}

In the immediate aftermath of the breakdown with Anthropic, OpenAI moved to fill the vacuum, securing its own \$200 million agreement to deploy its models on U.S. classified networks.\cite{guardianOpenAIMilitaryContract} The initial framework reportedly allowed the Pentagon to use OpenAI's systems for any \textit{lawful} purpose, while the company relied on internal technical safety constraints to uphold its stated red lines, such as prohibitions on directing autonomous weapons or enabling mass surveillance.\cite{nytOpenAIAmendsDeal} After a weekend of public backlash, OpenAI announced an amended contract that explicitly bans the intentional use of its technology for the mass surveillance of U.S. persons and clarifies that accountability for the use of force must remain with humans, thereby converging, at least at the level of written commitments, towards some of the safeguards Anthropic had been demanding.\cite{nytOpenAIAmendsDeal}

This triadic interaction between Anthropic, OpenAI, and the Pentagon operationalizes many of the safety concerns raised in purely simulated environments such as Rivera et al.'s escalation wargames or CosmoAgent's dark-forest scenarios (Section~\ref{sec:safety_analysis}). On the one hand, the Pentagon's insistence on unrestricted ``all lawful purposes'' access effectively pushes AI labs towards acting as providers of high-capability strategic agents whose action space includes real-world targeting, surveillance, and cyber-operations, a direct extension of the discrete and continuous action spaces studied in WarAgent, ALYMPICS, and Richelieu from the domain of \textit{simulation} to that of \textit{warfighting}.\cite{rivera2024escalation, mao2023alympics, guan2024richelieu} On the other hand, Anthropic's refusal to relax its usage policies, and OpenAI's subsequent attempt to thread a ``compromise,'' mirror exactly the design tensions discussed in Sections~\ref{sec:negotiation_protocols} and~\ref{sec:safety_analysis}: how far can safety guardrails, rules-oracle layers, and constitutional verifiers constrain powerful LLM agents before institutional actors perceive them as operationally unusable?

Ultimately, just as the deployment of Diella illustrates the integration of AI into civilian legislative and administrative processes (Section~\ref{sec:diella_real_world}), the Anthropic--OpenAI--Pentagon dispute marks the definitive entry of Large Language Models into the highest echelons of national security and military command. This transition from experimental wargaming (as explored by Rivera et al.) to live operational deployment represents a paradigm shift: AI is no longer just a tool for data analysis, but is rapidly becoming an active participant in the institutions that manage global crises and authorize the use of force.

As these models become embedded in the command-and-control infrastructure of global superpowers, a fundamental question emerges: how will these AI agents actually behave when tasked with navigating high-stakes geopolitical and military scenarios? The accelerated rate of this adoption leaves little room to empirically understand their strategic reasoning, propensity for escalation, or ideological biases under extreme pressure. Consequently, there is an urgent structural need within the research community for robust, deterministic testing methodologies. Before governments can safely rely on LLMs to interpret international threats or dictate military strategy, these systems must first proceed through rigorous auditing within causality-driven sandboxes capable of safely simulating the complex, multipolar realities they are now being deployed to govern.



\section{Summary of Limitations and Open Challenges}
\label{sec:soa_limitations}

While the current literature demonstrates the potential of LLMs in tactical gaming (e.g., \textit{Cicero}) and historical role-play (e.g., \textit{WarAgent}), a critical analysis reveals structural limitations that hinder the application of these systems to reliable, explainable geopolitical modeling. This thesis identifies four primary research gaps that define the scope of the proposed work.

\subsection*{1. The Explainability and Verifiability Gap}
The most significant limitation in current MAS is the lack of verifiable causal reasoning. While \textbf{WarAgent} provides unstructured text traces, these are often essentially "stream of consciousness" outputs that are difficult to query systematically.
Conversely, \textbf{Cicero} achieves high strategic performance but operates as a "black box" regarding its reasoning intent. Moreover, existing frameworks lack a quantitative bridge between internal reasoning and observable actions.
\\
\textbf{Gap:} There is currently no framework for formal, atomic transparency that records specific JSON-structured intents alongside a dedicated state-loading mechanism for counterfactual "What-If" verification.

\subsection*{2. Strategic Memorization and Topological Bias}
Historical plausibility in current simulations often relies on "Structural Fingerprinting."
Existing simulations are predominantly tied to the fixed, static topology of the Earth, creating a "Strategic Memorization" bias.
If an LLM recognizes the historical scenario through residual fingerprints, it may revert to reciting historical facts rather than simulating autonomous reasoning.
\\
\textbf{Gap:} Current architectures lack the ability to decouple strategic reasoning from geographic familiarity, preventing the assessment of whether agents truly comprehend geopolitical fundamentals (e.g., resource-density logistics or spatial choke points) independent of their pre-training data.

\subsection*{3. Fragmentation of Action and Validation Spaces}
The complexity of real-world diplomacy is poorly represented by current interaction protocols. Models typically rely on simplistic discrete moves or scalar bidding.
Furthermore, agent intentions are often decoupled from a rigid "physics engine," allowing for hallucinations of actions that violate logisitical or economic constraints.
\\
\textbf{Gap:} There is a lack of an integrated "Action Layer" that acts as a Rules Oracle, enforcing deterministic constraints (e.g., 15\% trade stock caps, pathfinding territory rules, and budget thresholds) before an agent's intent is admitted as reality.

\subsection*{4. The Decoupling of Social Dynamics and Production Physics}
Finally, there is a disconnect between top-down policy making and bottom-up physical capability.
In most models, "Public Opinion" is either a static variable or a purely descriptive element, rather than a deterministic constraint on the nation's industrial output.
\\
\textbf{Gap:} Existing frameworks lack a unified "Three Pillars" feedback loop where top-down government decisions (e.g., Welfare investment vs. War Taxes) dynamically and deterministically influence the nation's physical capability through a Satisfaction-Production engine.


\bigskip
\textbf{Proposed Contribution.} To address these limitations, this thesis introduces \textbf{GeoMAS}, a multi-agent framework situated in \textbf{synthetic environments} designed to decouple strategic agency from historical data leakage. The architecture establishes a \textbf{duality of determinism}, where a rigid physical-economic engine constrains and validates the \textbf{strategic intent} of LLM-powered agents in real-time. By integrating \textbf{causal tracing} and \textbf{divergent scenario analysis}, this work isolates the intrinsic reasoning capabilities of AI agents, providing a verifiable sandbox for the study of \textbf{geopolitical emergence} and the \textbf{social-material feedback loops} that govern national stability.

